% RESEARCH_DRAFT_SHA256: 7f14bd132670e6013291f0decec99b275978b4cf3b36ec24614a65dd4f6e4aa2
% Layout and source-reference wrappers expanded for the web renderer.
\documentclass[11pt]{article}
\usepackage[margin=25mm]{geometry}
\usepackage{fontspec,unicode-math}
\setmainfont{Latin Modern Roman}
\setmathfont{Latin Modern Math}
\usepackage{amsmath,amssymb,mathtools,xurl,hyperref}
\hypersetup{colorlinks=true,urlcolor=blue,linkcolor=blue}
\setcounter{secnumdepth}{0}
\setlength{\parindent}{0pt}
\setlength{\parskip}{5pt}
\setlength{\emergencystretch}{3em}
\newcommand{\sref}[2]{\hyperlink{source:#1:#2}{[S#2]}}
\newcommand{\cataloguescope}[1]{\paragraph{Source record.}#1}
\newcommand{\floor}[1]{\left\lfloor#1\right\rfloor}
\begin{document}
\section*{Research-attempt provenance}
This standalone research attempt retains the complete original definition body below,
including its historical source and preparation notes. Those notes describe the
earlier catalogue preparation; the new attempt and its verification are identified
separately after that body. No original ID, catalogue status, or source statement is changed.

\section{Sample Complexity of Active Coordinate Identification}\label{statement}

\subsection{Definitions and mathematical statement}

\paragraph{Shared notation.}Standard number systems are $\mathbb N,\mathbb Z,\mathbb Q,\mathbb R,\mathbb C$. Specialized notation and index ranges follow the cited source and the definitions below.


An active coordinate is a variable on which the target function depends. Nonadaptive sampling chooses all observation points before seeing any observed function values.
\paragraph{Record scope and qualifications.}
The catalogue title is retained as metadata. Named certificates, contracts, bounds or frameworks in that title are not asserted as results of the original source extract.
The smoothness class, approximation error, noise model, confidence, and meaning of approximately low-dimensional are source-specific.

Source-record limitation: The approximate low-dimensional function class, observation assumptions, and active-coordinate identification criterion are absent.
\paragraph{Statement recorded in the supplied source.}
In recovering an approximately low-dimensional function of $N$ variables from samples, how many non-adaptive sample points are required to identify the active coordinates? \href{https://doi.org/10.4171/owr/2010/33}{[S2]}
\paragraph{Source status.}
The supplied catalogue labels this source record open; that metadata was checked against the supplied snapshot on 2026-10-11. The label is preserved for provenance and is not an independent verification of present open status or mathematical correctness. The catalogue literature assessment is dated 2026-08-21.
\subsection{Short English statement}
Determine the sample count needed to identify variables used by an approximately low-dimensional function.
\subsection{Sources}
\begin{itemize}
\item[S1] Ulam AI and the UnsolvedMath Contributors, UnsolvedMath: A Curated Collection of Open Mathematics Problems, catalogue entry. Catalogue curation and metadata: CC BY 4.0. The source extract is retained; definitions, English summaries and qualifications are editorial additions.
\url{https://huggingface.co/datasets/ulamai/UnsolvedMath}.
\item[S2] Wavelet and Multiscale Methods (2011), Oberwolfach Reports, DOI 10.4171/owr/2010/33, pp. 1953--2025
\url{https://doi.org/10.4171/owr/2010/33}.
\item Source-specific Oberwolfach Reports licences, including CC BY-SA 4.0 where applicable, are retained; see THIRD\_PARTY\_NOTICES.md.
\end{itemize}
\paragraph{Preparation and review.}
The original source and UnsolvedMath contributors retain attribution for the source extract. Codex (GPT-6; exact serving revision not exposed) assisted with definitions, English summaries, scope notes and formatting. No new human mathematical authorship or independently verified solution is asserted. Catalogue identity checked against the complete supplied registry and current repository. Distinct upstream records, including repeated questions and broad or damaged extracts, are retained for complete coverage of the source CS category. A repeated record is not a new independent conjecture. Complete mathematical-equivalence and current-literature review has not been performed. Preparation and final review follow contribution prompts C and D, with source gaps disclosed. Structural and rendering checks do not establish a new theorem or open status.
\label{end:statement}

\section{Literature and context check}
The primary report on printed pages 2013--2014 explicitly warns that exact active-coordinate recovery may fail even in the exact functional model. It seeks an approximating function with appropriate coordinates, and states a nonadaptive approximation result of order $(L+1)^\ell\log N$ with dimension-dependent constants. This is more qualified than the catalogue's terse identification wording. The original text is retained; the attempt below tests exact support identification and gives a separate calibrated model where a sample bound can be proved.
\section{Research attempt: Finite-sample invisibility and a calibrated one-coordinate model}

\paragraph{Why finitely many point samples need not reveal a coordinate}
Fix arbitrary nonadaptive samples $x^{(1)},\ldots,x^{(m)}\in[0,1]^N$. For each coordinate $j$, form the nonconstant polynomial
\[
 q_j(t)=\prod_{r=1}^m(t-x_j^{(r)})^2
\]
when $m\geq1$. For $m=0$ take $q_j(t)=t$. It is not the zero polynomial, is nonconstant, and vanishes at every observed coordinate value in the sampled case. Thus $f_j(x)=a_jq_j(x_j)$ with $a_j\ne0$ genuinely depends only on coordinate $j$, yet all its observed sample values are zero. Distinct active-coordinate choices are indistinguishable. Scaling $a_j$ small enough also meets any fixed positive upper bound on a finite number of derivative sup norms, because the unscaled polynomial and its finitely many derivatives are bounded on $[0,1]$. No lower signal-strength condition is supplied by such smoothness upper bounds.

For $N\geq2$, any algorithm receiving this common zero data must fail exact identification on at least one of these functions. For a randomized estimator under a uniform choice among the $N$ functions, its success probability is at most $1/N$, because its same output distribution is used for all of them. This is a fixed-design obstruction; an adaptive deterministic algorithm can likewise be followed along its all-zero answer path and the final polynomials constructed at the resulting finitely many points. No claim about randomized adaptive designs with a function chosen beforehand is made.

In the approximate model, even the phrase ``the active coordinates'' can be nonunique: the zero function is within any positive tolerance $\varepsilon$ of many nonzero univariate functions of different coordinates with sup norm at most $\varepsilon$. Recovering an adequate approximating coordinate set is a different task from uniquely recovering a latent representation.

\paragraph{A sample count once the model is calibrated}
Assume instead $f(x)=a x_j$ for one unknown index $j$, where $a>0$. Query the all-one point first, and use $k=\lceil\log_2N\rceil$ additional Boolean points whose coordinate columns are distinct binary length-$k$ codes. The first value is $a$; subsequent values are either zero or $a$ and reveal the code of $j$. This proves a nonadaptive $k+1$ sample upper bound in this exact linear family.

If each sample has deterministic error at most $\eta$ and $a\geq4\eta>0$, compare each coded observation with half the first observation. For a zero code bit, its observation is at most $\eta$, while that threshold is at least $(a-\eta)/2>\eta$. For a one bit, its observation is at least $a-\eta$, while the threshold is at most $(a+\eta)/2<a-\eta$. Thus every bit is still recovered. The strict comparisons follow from $a>3\eta$.

A matching logarithmic information lower bound requires an access restriction. If the amplitude is fixed at one and samples must be Boolean points, there are at most $2^m$ possible answer strings for $m$ queries, so identifying all $N$ indices needs $m\geq\lceil\log_2N\rceil$. With arbitrary exact real samples, this lower bound is false: one point having pairwise distinct coordinate values identifies the index when $a=1$. Precision, noise and the admissible sample alphabet therefore materially affect any sample-complexity claim.

\section{Verification and adversarial audit}
Repeated sample-coordinate values do not invalidate the polynomial construction. Nonzero amplitude preserves genuine coordinate dependence even when made very small. The lower bound is explicitly fixed-design and distinguishes exact support from approximation. The positive-amplitude noisy coding argument checks both kinds of bit and the first sample's own noise. No finite-bit counting argument is applied to unbounded-precision real answers.
\paragraph{Audit corrections.} Recovered the primary source's warning about invisible variables; kept the auxiliary Boolean-query lower bound separate from arbitrary real sampling.
\paragraph{External computation.} None; symbolic arguments and boundary checks only.
\section{Scope and remaining gaps}
The broad smooth approximately low-dimensional function problem is not resolved by the calibrated linear model. The primary source already distinguishes exact support from adequate approximation; the present obstruction and explicit noisy coding route clarify which additional hypotheses are needed for an identification sample bound.
\paragraph{Integration review.} Codex (GPT-6; exact serving revision not exposed) reviewed the written arguments and expanded layout and source-reference wrappers. Repository metadata and the production web renderer were checked using Python 3.13 and JavaScript with MathJax 3.2.2. This model-assisted review is not independent human proof verification; no PDF compilation or formal proof checking is asserted.

\section{Completion Estimate}
\noindent\textbf{COMPLETION: 0\%}
The calculations below the preserved source statement investigate a concrete route,
but establish no substantiated new progress toward the entire original target.
Known special cases and failed sufficient conditions are not credited as a solution.
\section{References}
The original source references and catalogue attribution are preserved in the
Sources section above. The following URLs were actually consulted for this attempt;
retention of an additional original reference does not claim fresh consultation of it.
\begin{itemize}
\item \url{https://ems.press/content/serial-article-files/46289} (consulted 2026-10-11). 
\end{itemize}
\section{Model and contributor disclosure}
The parent Codex GPT-6 model has reviewed the written mathematical arguments at their declared scopes. This is model-assisted review, not independent human verification.
OpenAI gpt-6-astra, requested reasoning configuration ultra, generated this research
attempt on 2026-10-11. The exact serving revision is not exposed. Codex GPT-6
performed orchestration and repository-format review; its exact serving revision
is also not exposed. The mathematical route, boundary checks and final adversarial
self-audit followed contribution prompts A then B in
\texttt{docs/contribution\_prompts.md}, after reading \texttt{CONTRIBUTING.md}.
Python 3.14.2 assembled the standalone documents; any mathematical computation is
specified in Verification. No theorem prover was used. No human mathematical
authorship, substantial human proof correction, or independent human proof checking
is asserted. Model-assisted review, source retrieval, and structural checks do not
establish correctness or novelty. Full private reasoning transcripts are unavailable
as reproducibility artifacts; the complete arguments, inputs and outcomes actually
used are stated here or in the companion batch audit.
\section{Final status}
Unresolved as a new research contribution. This is a concrete route investigation,
not a claim of new complete resolution or a current-status certification. A known
answer reproduced for context does not become a new contribution. The catalogue's
source-status label is preserved; this attempt status does not assert that a
source-labelled solved question is open.
\end{document}
